% 3.模型假设.tex —— 三、模型假设
\section{模型假设}
本文采用以下假设，并区分其适用范围：
\begin{itemize}
  \item \textbf{几何与材料结构：}药材视为连续、均质且各向同性的介质。问题一至三在\textbf{固定几何}下求解、不计干燥收缩；问题四引入附件 2 的收缩半径 $R(t)$ 作为\textbf{已知移动边界}，收缩仅体现在径向半径、轴向长度与端部效应仍不建模（题目未给长度，沿用问题一至三口径）。温度与含水率仅沿径向变化，以长圆柱中段近似忽略轴向梯度。该近似不能说明端部效应已被验证，长时间全场结论限于径向模型。
  \item \textbf{物性（分问适用，需避免混用）：}问题一仅假设热物性为常数，采用附录 2 中的 $\rho,c_p,k$，水分扩散系数仍为 $D(C)$；问题二、三则改用附录 3 的变物性 $\rho(C),c_p(C),k(C),D(C,T)$；问题四改用\textbf{附录 4} 的变物性 $\rho(C),c_p(C),k(C),D(C,T)$（各系数与附录 3 不同）。\textbf{常热物性假设只适用于问题一，不适用于问题二、三、四}——后三问的热物性随含水率变化、扩散系数随含水率与温度变化，任何“物性为常数”的表述都不得延用到后三问。
  \item \textbf{初始状态：}初始温度和含水率在药材内部均匀，分别为 $T_0=28$ $^\circ$C、$C_0=2.55$ kg/kg（干基）。
  \item \textbf{边界与环境：}中心满足轴对称条件，表面换热和传质采用线性 Robin 边界，系数按题设取值。附件环境数据作分段线性插值；问题二、三在 $14400$ s（$4$ h）后保持最后记录值。水分边界将附件浓度作为有效平衡参照，其适用性取决于该传质关系的成立。
  \item \textbf{能量简化：}不计化学放热、水分生成和汽化潜热。问题二、三以 $\rho(C)c_p(C)\,\partial T/\partial t$ 表示局部热响应，未显式追踪组成变化及迁移水分带来的焓变化。因此不能把该热方程直接解释为完整变组分系统的焓守恒方程。
  \item \textbf{水分迁移：}水分服从有效 Fick 扩散关系。问题一的 $D$ 仅依赖 $C$，热质方程解耦；问题二、三的 $D$ 同时依赖 $C$ 和 $T$，并经热物性形成双向耦合；问题四在附录 4 的 $D(C,T)$ 之上，还经 Landau 变换引入移动边界带来的对流项 $+(E R'/R)\partial C/\partial E$。
\end{itemize}
